\subsection{Prime redistribution and globally fixed derivative directions}

\begin{theorem}[DP53: the cyclotomic redistribution preserves the leading bare-jet budget]\label{DP53}
Let $J=\operatorname{Jac}(y^2=x^7-1)$ with the good integral model of DP22, and let
\[
 \Phi=\{1,2,3\},\qquad\Psi=\{1,2,4\}
       \subset(\Z/7)^\times
\]
be respectively its actual covariant CM type and the comparison type induced from $\Q(\sqrt{-7})$.
For odd $p\ne7$, label the three actual invariant directions by the global integral differential basis
$\omega_i=x^{i-1}dx/y$, $i=1,2,3$.
Put
\[
 c_2(p)=\frac{p}{p^2-1},\qquad
 c_3(p)=\frac1{p-1}-\frac1{p^3-1},\qquad
 D(p)=(p-1)(p^3+1),
\]
\[
 T(p)=\frac{p(p^2+p+1)}{D(p)}.
\]
The sharp \emph{bare integral directional} slopes are
\[
\begin{array}{c|ccc|c}
p\bmod7&c_1^\Phi(p)&c_2^\Phi(p)&c_3^\Phi(p)&B_\Phi(p)\\\hline
1&0&0&0&0\\
2&p/(p^3-1)&p^2/(p^3-1)&c_3(p)&2c_3(p)\\
4&p^2/(p^3-1)&p/(p^3-1)&c_3(p)&2c_3(p)\\
3&p/D(p)&p^3/D(p)&p^2/D(p)&T(p)\\
5&p^3/D(p)&p/D(p)&p^2/D(p)&T(p)\\
6&c_2(p)&c_2(p)&c_2(p)&3c_2(p)
\end{array}
\]
Every selected direction of the induced type has slope
\[
 c^\Psi(p)=
 \begin{cases}0,&p\equiv1,2,4\pmod7,\\
 c_2(p),&p\equiv3,5,6\pmod7.\end{cases}
\]

For arbitrary nonnegative integers $n_1,n_2,n_3$, held fixed as the prime varies, let $N=n_1+n_2+n_3$. Uniformly in these orders and $X\ge2$,
\begin{align}
 \sum_{\substack{p\le X\\p\ne2,7}}\log p
       \sum_{i=1}^3 n_i c_i^\Phi(p)
       &=\frac N2\log X+O(N),\label{eq:cmprimeactual}\\
 \sum_{\substack{p\le X\\p\ne2,7}}\log p
       \sum_{i=1}^3 n_i c^\Psi(p)
       &=\frac N2\log X+O(N).\label{eq:cmprimeinduced}
\end{align}
Their difference is $O(N)$. The same bound holds for the full integer periodic-floor lower certificates of DP43, not merely their linear slopes.
These are statements about bare integral derivatives and their lower certificates, not upper bounds on actual theta-selected gain.
\end{theorem}

\begin{proof}
\emph{Actual integral directions at every good prime.}
The two-chart model of DP22 and its relative Jacobian are defined over $\Z[1/14]$. The displayed differentials form a basis on every good fibre, hence a primitive integral relative basis. Their dual Lie basis and the corresponding invariant directions $D_i$ are therefore rational and integral at every prime under consideration.
The covariant cyclotomic automorphism has cotangent and tangent characters $1,2,3$; the ordinary transpose retains each scalar eigenvalue.
After the unramified extension containing $\zeta_7$, the distinct seventh-root characters split the tangent lattice into rank-one direct summands. Thus a local Teichm\"uller tangent generator of a fixed character differs from this global generator only by a unit. No prime-dependent direction is being selected.

Let $f=\operatorname{ord}_7(p)$. The CM height/type theorem gives height $f$ to each Frobenius orbit and dimension equal to its intersection with $\Phi$; this is the standard local formula already used in DP22, \cite[\S5, (10*) and Theorem 3, pp.~106--108]{Tate352}.
An orbit with positive selected count is connected: a nonzero \'etale quotient would carry the unramified degree-$f$ field action, so its rational Tate dimension would be a positive multiple of $f$. It would have full height and force dimension zero, a contradiction. For a partially selected orbit the dual is connected by the same argument.
The integral idempotents and the connected formal-group equivalence identify actual formal factors, not merely isogeny factors.

For each partially selected orbit, apply the integral argument of DP45 and DP49. Over the $p$-nilpotent quotients of the unramified base, its display has rank $f$; the Witt character idempotents are free summands, semilinear Frobenius permutes them cyclically, and its linearization becomes invertible after inverting $p$. The summands therefore all have rank one. The exact Hodge quotient forces $Q_i=IP_i$ on selected characters and $Q_i=P_i$ on the others. The normal-decomposition matrix is monomial and invertible, so all arrows are units. Compatible generators lift through the Witt reductions. This gives the actual integral display required by DP40 at every such prime.

For the actual type, the selected full orbits occur only for $p\equiv1$, when $f=1$. Each is a height-one formal group with sharp slope zero. To check exactness directly, kill logarithm terms of degrees $2,\ldots,p-1$ by an integral strict coordinate. If $b_p$ is the next coefficient, height one makes $pb_p$ a unit, since the degree-$p$ coefficient of $[p]$ is $(p-p^p)b_p$.
The integral translation coefficient is $D^p/p!+b_pD$.
Hence $D^p$ reduces to a unit multiple of $D$, and every iterated order $p^r$ takes a unit value on a normalized coordinate. Integrality gives the opposite slope bound zero.
For the induced comparison type, the split factors are ordinary of multiplicative type. After a faithfully unramified strictly Henselian extension their \'etale Cartier duals are constant and their formal groups are multiplicative. The restricted operation is invertible; a CM eigendirection has unit $p^f$-power, again giving sharp slope zero. In the other classes the type alternates, so every selected gap is two and DP40 gives $c_2(p)$ in each direction.

\emph{The globally labelled spectrum.}
For $a=p\bmod7$ and $i\in\Phi$, the exact cyclic formula of DP40--DP43 is
\begin{equation}\label{eq:cmprimecycle}
 c_i^\Phi(p)
 =\sum_{r\ge1}\mathbf1_{a^ri\notin\Phi}p^{-r}
 =\frac{\sum_{r=1}^f\mathbf1_{a^ri\notin\Phi}p^{f-r}}{p^f-1}.
\end{equation}
It also gives zero on the selected ordinary orbits.
For $a=2$, use $(1,2,4)$ and $(3,6,5)$; for $a=4$, use their inverse orders. These give the two height-three rows of the table.
For $a=3$, the order $(1,3,2,6,4,5)$ has its first three positions selected; for $a=5$, use $(2,3,1,5,4,6)$. DP45's ordered triple $(p,p^2,p^3)/D(p)$ therefore becomes respectively $(p,p^3,p^2)/D(p)$ and $(p^3,p,p^2)/D(p)$ in the fixed global label order.
For $a=6$, the orbits $(1,6),(2,5),(3,4)$ each contain one selected index. This proves the whole table. The induced type is constant on its orbits for $a\in\Psi$ and alternating for $a\notin\Psi$, proving its formula as well.

As an exact common-variable identity for $p>1$,
\begin{equation}\label{eq:cmprimeformalaverage}
 \frac{4c_3(p)+2T(p)+3c_2(p)}6-\frac{9c_2(p)}6
       =\frac{4p^2(p^2+1)}{3(p^6-1)}.
\end{equation}
Indeed, over denominator $p^6-1$, the numerators of $c_3,T,c_2$ are
$p(p+1)(p^3+1)$, $p(p^2+p+1)^2$, and $p(p^4+p^2+1)$.
The numerator of $4c_3+2T-6c_2$ is $8p^2(p^2+1)$.
This formal averaging identity does not determine the sign of the actual prime-sum constant.

\emph{The unconditional prime sum.}
Set
\[
 \alpha_i(a)=\mathbf1_{ai\notin\Phi},\qquad
 \alpha^\Psi(a)=\mathbf1_{a\in\{3,5,6\}}.
\]
The cyclic series gives the uniform tail estimate
\[
 0\le c_i^\Phi(p)-\alpha_i(a)/p
       \le\sum_{r\ge2}p^{-r}=\frac1{p(p-1)},
\]
and the same estimate holds for the induced type.
For each fixed $i$, multiplication by $i$ permutes the six nonzero residues, exactly three of which land outside $\Phi$. Therefore
\[
 \sum_{a=1}^6\alpha_i(a)=3
       =\sum_{a=1}^6\alpha^\Psi(a).
\]
Use the established fixed-modulus prime-sum estimate
\begin{equation}\label{eq:cmprimeAP}
 \sum_{\substack{p\le X\\p\equiv a\ (7)}}\frac{\log p}{p}
       =\frac16\log X+O(1),\qquad a=1,\ldots,6.
\end{equation}
An exact direct reference is Rudnick, \emph{Primes in Arithmetic Progressions: Fixed Modulus}, May 26, 2015, \S2, p.~2, equation~(1); its preceding paragraph gives the unconditional quantitative fixed-modulus PNT.\footnote{Official author PDF: \url{https://www.math.tau.ac.il/~rudnick/courses/sieves2015/fixed\%20progressions.pdf}. The original page and fixed-modulus hypotheses were checked. The known PNT is cited, not reproved.}
The tail errors weighted by $\log p$ converge absolutely, by comparison with $\sum_{m\ge2}(\log m)/m^2$. Thus for each globally fixed direction
\[
 \sum_{\substack{p\le X\\p\ne2,7}}c_i^\Phi(p)\log p
       =\frac12\log X+O(1),\qquad
 \sum_{\substack{p\le X\\p\ne2,7}}c^\Psi(p)\log p
       =\frac12\log X+O(1).
\]
The finitely many bad primes affect only the bounded terms.
The three labelwise errors are uniformly bounded for all $X\ge2$.
Multiplying by $n_i$ and adding proves
\eqref{eq:cmprimeactual}--\eqref{eq:cmprimeinduced}, with constants independent of the orders. These orders may vary with an outer parameter; they must be fixed while the prime in the sum varies.

\emph{The integer periodic-floor certificates.}
Write
\[
 E_{i,p}^\Phi(n)
       =\sum_{r\ge1}\mathbf1_{a^ri\notin\Phi}
          \floor{n/p^r}.
\]
The induced certificate $E_p^\Psi(n)$ is zero in the split classes and
$\sum_{r\text{ odd}}\floor{n/p^r}$ in the inert classes.
These are the actual coefficientwise lower certificates of DP43; pointwise exactness is not asserted.
We prove, uniformly in $X$,
\begin{equation}\label{eq:cmprimefloors}
 \left|\sum_{\substack{p\le X\\p\ne2,7}}\log p
       \sum_{i=1}^3\bigl(E_{i,p}^\Phi(n_i)-E_p^\Psi(n_i)\bigr)\right|
       =O(N).
\end{equation}
Fix $i$ and put $n=n_i$. If $n\le1$, all floors vanish.
Otherwise put $Y=\min(X,n)$ and $\delta_i(a)=\alpha_i(a)-\alpha^\Psi(a)$.
The first-power contribution is
\begin{align*}
 &\sum_{\substack{p\le Y\\p\ne2,7}}
       \delta_i(a)\floor{n/p}\log p\\
 &\qquad=n\sum_{\substack{p\le Y\\p\ne2,7}}
       \delta_i(a)\frac{\log p}{p}
       +O\left(\sum_{p\le Y}\log p\right)=O(n).
\end{align*}
The first sum is bounded by \eqref{eq:cmprimeAP} because $\sum_a\delta_i(a)=0$; the last sum is $O(Y)=O(n)$ by the same fixed-modulus PNT summed over the six classes.
The absolute higher-power difference is at most
\[
 n\sum_p\log p\sum_{r\ge2}p^{-r}
       =n\sum_p\frac{\log p}{p(p-1)}=O(n).
\]
Adding over the three labels proves \eqref{eq:cmprimefloors}.
This checks rounding and the entire periodic pattern, not only its first real slope.
\end{proof}

For equal direction orders the actual first coefficients across the six classes are $(0,2,1,2,1,3)$, whereas the induced type gives $(0,0,3,0,3,3)$. Their sums are both nine. Some prime classes gain while others lose; the signed first-order terms require prime equidistribution, not an absolute same-prime summability argument.

The combinatorial reason has a precise general form. If a finite group $\Gamma$ acts transitively on $\Omega$ with $|\Omega|=2g$ and $|A|=g$, then
\[
 \frac1{|\Gamma|}\sum_{\sigma\in\Gamma}
       \#\{i\in A:\sigma i\notin A\}=\frac g2.
\]
For each fixed $i$, half the equally frequent targets lie outside $A$; fixed nonnegative weights $n_i$ give average $\tfrac12\sum_i n_i$ by the same count.
In the present abelian cyclotomic case, Frobenius is the actual residue element and \eqref{eq:cmprimeAP} supplies this distribution. For a nonabelian embedding action the displayed function need not be constant on conjugacy classes, and sampling primes of a reflex or other base field can change the distribution. The group average alone is not asserted to prove a general prime-average theorem for every CM variety.

Thus the concrete Jacobian does redistribute bare derivative divisibility into prime classes where the induced type has none, but for arbitrary orders assigned to its fixed global rational directions the linear and periodic-floor budgets change by only $O(N)$. Their leading logarithmic coefficient is unchanged. This does not identify them with normalized theta-selected values, reverse the source's one-way \path{LinkG} lower certificate, or permit prime-dependent reassignment of prescribed directions. A finite $O(N)$ improvement can still matter in a finite inequality. Special sections, selected harmonic weights, the new endpoint period and the actual normalization/cost terms remain additional data.

